\subsection{Cálculos numéricos}

\ejemplosboxed[10]{Realiza las siguientes operaciones de \textit{cálculo numérico}:

		\begin{parts}
			\begin{multicols}{2}

				% \addcontentsline{toc}{subsection}{Suma de números}
				\subsubsection*{Suma de números}
				
			%	\part $849.332+242.25+469.381=$ \fillin[$1560.963$][0in]
				
				%\part $27.05+34.99+0.1=$ \fillin[$62.14$][0in]
				
				\part $0.1+0.02+0.03+0.4=$ \fillin[$0.55$][0in]
				
			%	\part $0.11+2+3.8=$ \fillin[$5.91$][0in]

				% \addcontentsline{toc}{subsection}{Resta de números}
				\subsubsection*{Resta de números}
				
			%	\part $4934-451-682=$ \fillin[$3801$][0in]
				
			%	\part $0.1-0.02=$ \fillin[$0.08$][0in]
				
				\part $0.1-0.02-0.03-0.4=$ \fillin[$-0.35$][0in]
				
			%	\part $0.11-2-3.8=$ \fillin[$-5.69$][0in]

				% \addcontentsline{toc}{subsection}{Multiplicación de números}
				\subsubsection*{Multiplicación de números}
				
				%\part $19.3\times6.27=$ \fillin[$121.011$][0in]
				
				\part $0.1\times0.02=$ \fillin[$0.002$][0in]
				
				%\part $100.1\times 0.99=$ \fillin[$99.099$][0in]
				
				%\part $0.11\times2\times3.8=$ \fillin[$0.836$][0in]

				% \addcontentsline{toc}{subsection}{División de números}
				\subsubsection*{División de números}
				
				\part $922\divisionsymbol1.2=$ \fillin[$768.333$][0in]
				
			%	\part $0.1\divisionsymbol0.02=$ \fillin[$5$][0in]
				
			%	\part $180\divisionsymbol 0.09=$ \fillin[$2000$][0in]
				
			%	\part $25.25\divisionsymbol 0.5=$ \fillin[$50.5$][0in]

				% \addcontentsline{toc}{subsection}{Resolución de problemas}
				\subsubsection*{Resolución de problemas}
				\part Una computadora tiene un disco duro de 368 GB de memoria, si varios programas ocupan 128.75 GB. ¿Qué cantidad de memoria está libre?

				\begin{solutionbox}{1.2cm}
					$368-128.75=$ \fillin[$239.25$][0in]
				\end{solutionbox}
				
			\end{multicols}
		\end{parts}}

\questionboxed[10]{Realiza las siguientes operaciones de \textit{cálculo numérico}:

		\begin{parts}
			\begin{multicols}{2}

				% \addcontentsline{toc}{subsection}{Suma de números}
				\subsubsection*{Suma de números}
				
				\part $849.332+242.25+469.381=$ \fillin[$1560.963$][0in]
				
				\part $27.05+34.99+0.1=$ \fillin[$62.14$][0in]
				
				\part $0.1+0.02+0.03+0.4=$ \fillin[$0.55$][0in]
				
				\part $0.11+2+3.8=$ \fillin[$5.91$][0in]

				% \addcontentsline{toc}{subsection}{Resta de números}
				\subsubsection*{Resta de números}
				
				\part $4934-451-682=$ \fillin[$3801$][0in]
				
				\part $0.1-0.02=$ \fillin[$0.08$][0in]
				
				\part $0.1-0.02-0.03-0.4=$ \fillin[$-0.35$][0in]
				
				\part $0.11-2-3.8=$ \fillin[$-5.69$][0in]

				% \addcontentsline{toc}{subsection}{Multiplicación de números}
				\subsubsection*{Multiplicación de números}
				
				\part $19.3\times6.27=$ \fillin[$121.011$][0in]
				
				\part $0.1\times0.02=$ \fillin[$0.002$][0in]
				
				\part $100.1\times 0.99=$ \fillin[$99.099$][0in]
				
				\part $0.11\times2\times3.8=$ \fillin[$0.836$][0in]

				% \addcontentsline{toc}{subsection}{División de números}
				\subsubsection*{División de números}
				
				\part $922\divisionsymbol1.2=$ \fillin[$768.333$][0in]
				
				\part $0.1\divisionsymbol0.02=$ \fillin[$5$][0in]
				
				\part $180\divisionsymbol 0.09=$ \fillin[$2000$][0in]
				
				\part $25.25\divisionsymbol 0.5=$ \fillin[$50.5$][0in]

				% \addcontentsline{toc}{subsection}{Resolución de problemas}
				\subsubsection*{Resolución de problemas}
				
				
                \part Un elástico se estira tres veces su longitud en su estado normal. Si mide 5.23 cm en su estado normal, ¿cuántos centímetros alcanza al ser estirado?
			
				\begin{solutionbox}{1.2cm}
					$5.23 \times 3=$ \fillin[$15.69$][0in]
				\end{solutionbox}
				
				% \part Entre José y su hermano están arreglando el jardín de su casa. José arregló $\dfrac{3}{8}$ del jardín y su hermano $\dfrac{1}{4}$. ¿Qué parte del jardín han arreglado?

				% \begin{solutionbox}{3cm}
				% 	\[\dfrac{3}{8}+\dfrac{1}{4}=\dfrac{3}{8}+\dfrac{2}{8}=\dfrac{5}{8}\]
				% \end{solutionbox}
			\end{multicols}
		\end{parts}}

    % \subsubsection{Suma de números}
	% \questionboxed[2]{Realiza las siguientes \emph{sumas}:

	% 	\begin{parts}
	% 		\begin{multicols}{3}
	% 			\part \ifprintanswers{\opadd[hfactor=decimal,resultstyle=\color{red},carryadd=true]{464}{303}}
	% 			\else{\opadd[hfactor=decimal,resultstyle=\color{white},carryadd=false]{464}{303}}\fi

	% 			\part \ifprintanswers{\opadd[hfactor=decimal,resultstyle=\color{red},carryadd=true]{647.94}{564.973}}
	% 			\else{\opadd[hfactor=decimal,resultstyle=\color{white},carryadd=false]{647.94}{564.973}}\fi

	% 			\part \ifprintanswers{\opadd[hfactor=decimal,resultstyle=\color{red},carryadd=true]{98.97}{46.52}}
	% 			\else{\opadd[hfactor=decimal,resultstyle=\color{white},carryadd=false]{98.97}{46.52}}\fi

	% 			\columnbreak%

	% 			\part \ifprintanswers{\opadd[hfactor=decimal,resultstyle=\color{red},carryadd=true]{5423}{3214}}
	% 			\else{\opadd[hfactor=decimal,resultstyle=\color{white},carryadd=false]{5423}{3214}}\fi

	% 			\part \ifprintanswers{\opadd[hfactor=decimal,resultstyle=\color{red},carryadd=true]{344.64}{280.79}}
	% 			\else{\opadd[hfactor=decimal,resultstyle=\color{white},carryadd=false]{344.64}{280.79}}\fi

	% 			\part \ifprintanswers{\opadd[hfactor=decimal,resultstyle=\color{red},carryadd=true]{67.67}{52.97}}
	% 			\else{\opadd[hfactor=decimal,resultstyle=\color{white},carryadd=false]{67.67}{52.973}}\fi

	% 			\columnbreak%

	% 			\part \ifprintanswers{\opadd[hfactor=decimal,resultstyle=\color{red},carryadd=true]{54.54}{19.23}}
	% 			\else{\opadd[hfactor=decimal,resultstyle=\color{white},carryadd=false]{54.54}{19.23}}\fi

	% 			\part $321+51+134=$ \fillin[$506$][0in]

	% 			\part $0.1+0.02+0.03+0.4=$ \fillin[$0.55$][0in]
	% 		\end{multicols}
	% 	\end{parts}
	% }

	% \subsubsection{Resta de números}

	% \questionboxed[2]{Realiza las siguientes \emph{restas}:

	% 	\begin{parts}
	% 		\begin{multicols}{3}
	% 			\part \ifprintanswers{   \opsub[hfactor=decimal,resultstyle=\color{red},carryadd=true,carrysub=true]{82.48}{28.19} }
	% 			\else{            \opsub[hfactor=decimal,resultstyle=\color{white},carryadd=false,carrysub=false]{82.48}{28.19}}\fi

	% 			\part \ifprintanswers{   \opsub[hfactor=decimal,resultstyle=\color{red},carryadd=true,carrysub=true]{495}{92} }
	% 			\else{            \opsub[hfactor=decimal,resultstyle=\color{white},carryadd=false,carrysub=false]{495}{92}}\fi

	% 			\part \ifprintanswers{   \opsub[hfactor=decimal,resultstyle=\color{red},carryadd=true,carrysub=true]{937}{682} }
	% 			\else{            \opsub[hfactor=decimal,resultstyle=\color{white},carryadd=false,carrysub=false]{937}{682}}\fi

	% 			\part \ifprintanswers{   \opsub[hfactor=decimal,resultstyle=\color{red},carryadd=true,carrysub=true]{341}{212} }
	% 			\else{            \opsub[hfactor=decimal,resultstyle=\color{white},carryadd=false,carrysub=false]{341}{212}}\fi

	% 			\part \ifprintanswers{   \opsub[hfactor=decimal,resultstyle=\color{red},carryadd=true,carrysub=true]{465.76}{292.41} }
	% 			\else{            \opsub[hfactor=decimal,resultstyle=\color{white},carryadd=false,carrysub=false]{465.76}{292.41}}\fi

	% 			\part \ifprintanswers{   \opsub[hfactor=decimal,resultstyle=\color{red},carryadd=true,carrysub=true]{762}{394} }
	% 			\else{            \opsub[hfactor=decimal,resultstyle=\color{white},carryadd=false,carrysub=false]{762}{394}}\fi

	% 			\part \ifprintanswers{   \opsub[hfactor=decimal,resultstyle=\color{red},carryadd=true,carrysub=true]{850}{472} }
	% 			\else{            \opsub[hfactor=decimal,resultstyle=\color{white},carryadd=false,carrysub=false]{850}{472}}\fi

	% 			\part \ifprintanswers{   \opsub[hfactor=decimal,resultstyle=\color{red},carryadd=true,carrysub=true]{945}{173} }
	% 			\else{            \opsub[hfactor=decimal,resultstyle=\color{white},carryadd=false,carrysub=false]{945}{173}}\fi

	% 			\part \ifprintanswers{   \opsub[hfactor=decimal,resultstyle=\color{red},carryadd=true,carrysub=true]{0.1}{0.02} }
	% 			\else{            \opsub[hfactor=decimal,resultstyle=\color{white},carryadd=false,carrysub=false]{0.1}{0.02}}\fi
	% 		\end{multicols}
	% 	\end{parts}
	% }

	% \subsubsection{Multiplicación de números}

	% \questionboxed[2]{Realiza las siguientes \emph{multiplicaciones}:

	% 	\begin{parts}
	% 		\begin{multicols}{4}
	% 			\part \ifprintanswers{\opmul[hfactor=decimal,resultstyle=\color{red},displayintermediary=None]{284}{31} }
	% 			\else{\opmul[hfactor=decimal,resultstyle=\color{white},displayintermediary=None]{284}{31}}\fi\\[1em]

	% 			\part \ifprintanswers{\opmul[hfactor=decimal,resultstyle=\color{red},displayintermediary=None]{26.37}{13} }
	% 			\else{\opmul[hfactor=decimal,resultstyle=\color{white},displayintermediary=None]{26.37}{13}}\fi\\[1em]

	% 			\part \ifprintanswers{\opmul[hfactor=decimal,resultstyle=\color{red},displayintermediary=None]{411}{4} }
	% 			\else{\opmul[hfactor=decimal,resultstyle=\color{white},displayintermediary=None]{411}{4}}\fi\\[1em]

	% 			\part \ifprintanswers{\opmul[hfactor=decimal,resultstyle=\color{red},displayintermediary=None]{57}{1.39} }
	% 			\else{\opmul[hfactor=decimal,resultstyle=\color{white},displayintermediary=None]{57}{1.39}}\fi\\[1em]

	% 			\part \ifprintanswers{\opmul[hfactor=decimal,resultstyle=\color{red},displayintermediary=None]{24.13}{15} }
	% 			\else{\opmul[hfactor=decimal,resultstyle=\color{white},displayintermediary=None]{24.13}{15}}\fi\\[1em]

	% 			\part \ifprintanswers{\opmul[hfactor=decimal,resultstyle=\color{red},displayintermediary=None]{1851}{21} }
	% 			\else{\opmul[hfactor=decimal,resultstyle=\color{white},displayintermediary=None]{1851}{21}}\fi\\[1em]

	% 			\part \ifprintanswers{\opmul[hfactor=decimal,resultstyle=\color{red},displayintermediary=None]{515}{37} }
	% 			\else{\opmul[hfactor=decimal,resultstyle=\color{white},displayintermediary=None]{515}{37}}\fi\\[1em]

	% 			\part \ifprintanswers{\opmul[hfactor=decimal,resultstyle=\color{red},displayintermediary=None]{225}{9} }
	% 			\else{\opmul[hfactor=decimal,resultstyle=\color{white},displayintermediary=None]{225}{9}}\fi\\[1em]
	% 		\end{multicols}
	% 	\end{parts}
	% }

	% \subsubsection{División de números}

	% \questionboxed[2]{Calcula el \textbf{cociente} y el \textbf{residuo} de las siguientes \emph{divisiones}:

	% 	\begin{parts}
	% 		\begin{multicols}{3}
	% 			\part $785 \divisionsymbol 125=$ \\[1em]
	% 			Cociente: \fillin[6][0in]\\
	% 			Residuo: \fillin[35][0in]\\

	% 			\part $655.23 \divisionsymbol 23=$ \\[1em]
	% 			Cociente: \fillin[28][0in]\\
	% 			Residuo: \fillin[11][0in]\\

	% 			\part $123 \divisionsymbol 1.2=$ \\[1em]
	% 			Cociente: \fillin[102][0in] \\
	% 			Residuo: \fillin[6][0in]\\

	% 			\part $723 \divisionsymbol 8=$ \\[1em]
	% 			Cociente: \fillin[90][0in]\\
	% 			Residuo: \fillin[3][0in]\\

	% 			\part $90 \divisionsymbol 21=$ \\[1em]
	% 			Cociente: \fillin[4][0in]\\
	% 			Residuo: \fillin[6][0in]\\

	% 			\part $22 \divisionsymbol 0.2=$ \\[1em]
	% 			Cociente: \fillin[110][0in]\\
	% 			Residuo: \fillin[0][0in]\\
	% 		\end{multicols}
	% 	\end{parts}
	% }

	% \subsubsection{Resolución de problemas}

	% \questionboxed[2]{Resuelve los siguientes problemas:

	% 	\begin{parts}
	% 		\begin{multicols}{2}
	% 			\part Una computadora tiene un disco duro de 368 GB de memoria, si varios programas ocupan 128.75 GB. ¿Qué cantidad de memoria está libre?

	% 			\begin{solutionbox}{1.2cm}
	% 				$368-128.75=$ \fillin[$239.25$][0in]
	% 			\end{solutionbox}

	% 			\part En un estacionamiento conté 57 automóviles, 31 camionetas y 23 taxis, ¿cuántos vehículos había en total?

	% 			\begin{solutionbox}{1.2cm}
	% 				$57+31+23=$ \fillin[$111$][0in]
	% 			\end{solutionbox}

	% 			\part El precio de 385 artículos comerciales es de 1232 pesos. ¿Cuál es el precio unitario de cada artículo?

	% 			\begin{solutionbox}{1.2cm}
	% 				$1232 \divisionsymbol 385=$ \fillin[$3.20$][0in]
	% 			\end{solutionbox}
				
	% 			\part Las ventas de boletos que registra un cine en un fin de semana son las siguientes: 490 boletos vendidos el viernes, 780 el sábado y 1234 el domingo. ¿Cuántos boletos se vendieron en total?
				
	% 			\begin{solutionbox}{1.2cm}
	% 				$490+780+1234=$ \fillin[$2504$][0in]
	% 			\end{solutionbox}

	% 			\part Una pintura tiene un costo de 25.75 pesos el litro, una persona compra 48 litros. ¿Cuánto debe pagar?
			
	% 			\begin{solutionbox}{1.2cm}
	% 				$25.75 \times 48=$ \fillin[$1236$][0in]
	% 			\end{solutionbox}
			
	% 			\part Un elástico se estira tres veces su longitud en su estado normal. Si mide 5.23 cm en su estado normal, ¿cuántos centímetros alcanza al ser estirado?
			
	% 			\begin{solutionbox}{1.2cm}
	% 				$5.23 \times 3=$ \fillin[$15.69$][0in]
	% 			\end{solutionbox}

	% 			\part Si un dólar equivale a 19 pesos. ¿Cuántos pesos equivaldrán 615 dólares?
			
	% 			\begin{solutionbox}{1.2cm}
	% 				$19 \times 615=$ \fillin[$11685$][0in]
	% 			\end{solutionbox}

	% 			\part En un recipiente con agua, se agregaron otros 12.56 litros,llegando a completar 15.89 litros de agua. ¿Cuántos litros de agua había inicialmente en el recipiente?
			
	% 			\begin{solutionbox}{1.2cm}
	% 				$15.89 \times 12.56=$ \fillin[$3.33$][0in]
	% 			\end{solutionbox}
	% 		\end{multicols}
	% 	\end{parts}
	% }
   % \newpage
